\subsection{EXAMPLE 1.4: FRM EXAM 2009---QUESTION 1-6 (Jensen's alpha)}

\textbf{Enunciado}

\emph{Suppose Portfolio A has an expected return of 8\%, volatility of 20\%, and beta of 0.5. Suppose the market has an expected return of 10\% and volatility of 25\%. Finally, suppose the risk-free rate is 5\%. What is Jensen’s alpha for Portfolio A?}
\begin{enumerate}[label=\alph*.]
\item 10.0\%
\item 1.0\%
\item \hl{0.5\%}
\item 15\%
\end{enumerate}

\subsubsection*{Solución}

La medida de \emph{Jensen's alpha} compara el rendimiento esperado del portafolio
con el rendimiento que \emph{debería} tener según el CAPM. Su fórmula es

\[
\alpha_J
=
\mathbb E[R_A]
-
\Bigl(
R_f+\beta_A(\mathbb E[R_M]-R_f)
\Bigr).
\]

Aquí:

\begin{itemize}
    \item \(\mathbb E[R_A]\) es el rendimiento esperado del portafolio \(A\),
    \item \(R_f\) es la tasa libre de riesgo,
    \item \(\mathbb E[R_M]\) es el rendimiento esperado del mercado,
    \item \(\beta_A\) mide la exposición del portafolio al riesgo sistemático del mercado.
\end{itemize}

\paragraph{Paso 1. Prima de riesgo del mercado}

La prima de riesgo del mercado es

\[
\mathbb E[R_M]-R_f=0.10-0.05=0.05.
\]

Es decir, el mercado ofrece un rendimiento adicional de \(5\%\) sobre la tasa libre de riesgo.

\paragraph{Paso 2. Prima de riesgo correspondiente al portafolio}

Como el portafolio tiene \(\beta_A=0.5\), su exposición al riesgo sistemático
es la mitad de la del mercado. Por tanto, la prima de riesgo que le corresponde es

\[
\beta_A(\mathbb E[R_M]-R_f)=0.5(0.05)=0.025.
\]

\paragraph{Paso 3. Rendimiento requerido por CAPM}

El rendimiento requerido por CAPM para el portafolio \(A\) es

\[
R_f+\beta_A(\mathbb E[R_M]-R_f)
=
0.05+0.025
=
0.075.
\]

Es decir, un portafolio con \(\beta=0.5\) debería ofrecer un rendimiento esperado de \(7.5\%\).

\paragraph{Paso 4. Calcular el alpha}

Dado que el rendimiento esperado observado del portafolio es

\[
\mathbb E[R_A]=0.08,
\]

entonces

\[
\alpha_J
=
0.08-0.075
=
0.005.
\]

En porcentaje,

\[
\alpha_J=0.5\%.
\]


El portafolio ofrece \(0.5\%\) más de rendimiento del que el CAPM predice para su nivel
de riesgo sistemático. Por ello, su Jensen's alpha es positivo.









